% !TeX root = ../main.tex
% ---------------------------------------------------------------------------
% Risk Management, Contracts and Project Delivery
% Scope extension: not derived from the RSE3010 weekly material.
% Key sources: ITIG Code of Practice for Risk Management of Tunnel Works;
% ITA WG2 Guidelines for Tunnelling Risk Management; Essex, Geotechnical
% Baseline Reports for Construction; NEC and FIDIC forms.
% ---------------------------------------------------------------------------
\chapter{Risk Management, Contracts and Project Delivery}
\label{ch:risk-delivery}


\section{The nature of ground risk}
\label{sec:the-nature-of-ground-risk}

\subsection{The nature of risk in underground construction}
\label{sec:underground-risk}
\scopenote{Why ground risk is different: it cannot be eliminated by
investigation, it is unevenly distributed along an alignment, and the party
best placed to manage it is not always the party carrying it.}

\subsection{Risk identification and registers}
\label{sec:risk-registers}
\scopenote{Structured hazard identification, register format, ownership and
live maintenance through design and construction; the register as a working
document rather than a compliance artefact.}

\subsection{Qualitative and quantitative risk assessment}
\label{sec:risk-assessment}
\scopenote{Likelihood-consequence matrices, their known weaknesses, and
quantitative schedule and cost risk analysis where the exposure justifies it.}

\section{Codes of practice and baselines}
\label{sec:codes-of-practice-and-baselines}

\subsection{The ITIG Code of Practice}
\label{sec:itig}
\scopenote{The insurance-driven code that reshaped tunnelling risk practice
after a sequence of major losses, and the obligations it places on client,
designer and contractor.}

\subsection{Geotechnical Baseline Reports}
\label{sec:gbr}
\scopenote{Baselining ground conditions contractually: what a baseline is and
is not, how it differs from a factual and an interpretative report, and the
discipline required to write one that is actually testable.}

\subsection{Differing site conditions and claims}
\label{sec:differing-site-conditions}
\scopenote{The mechanism by which encountered conditions are compared against
the baseline, and how disputes typically arise from ambiguous baselining.}

\section{Contracts and procurement}
\label{sec:contracts-and-procurement}

\subsection{Contract models and risk allocation}
\label{sec:contract-models}
\scopenote{Design-bid-build, design and construct, early contractor
involvement, alliance and collaborative models; NEC and FIDIC forms; and how
each allocates ground risk.}

\subsection{Procurement and tendering for underground works}
\label{sec:procurement}
\scopenote{Prequalification, reference design versus performance
specification, tender period investigation, and the false economy of
transferring unpriceable risk.}

\subsection{Dispute avoidance and resolution}
\label{sec:disputes}
\scopenote{Dispute review boards and dispute adjudication boards, early
warning mechanisms, and the record-keeping that makes resolution possible.}

\section{Cost and programme}
\label{sec:cost-and-programme}

\subsection{Cost estimation and contingency}
\label{sec:cost-contingency}
\scopenote{Rate build-up for underground works, the influence of ground class
distribution, reference class forecasting, and probabilistic contingency
setting.}

\subsection{Programme, productivity and schedule risk}
\label{sec:programme}
\scopenote{Advance rates by method and ground class, utilisation, learning
curve, and the schedule consequences of interventions and stoppages.}

\section{Safety and stakeholders}
\label{sec:safety-and-stakeholders}

\subsection{Health and safety management}
\label{sec:safety-management}
\scopenote{Safety in design obligations, construction safety management
systems, competence and training, and the Australian regulatory framework for
tunnelling.}

\subsection{Stakeholder, community and environmental approvals}
\label{sec:stakeholders}
\scopenote{EES and EIS processes, community engagement, noise vibration and
settlement complaints, and the reputational risk of poor communication.
Framework introduced in Chapter~\ref{ch:context}.}
